\title{Group Sequence Policy Optimization}

\begin{abstract}
We present Group Sequence Policy Optimization (GSPO), a reinforcement learning algorithm tailored for training large language models, which emphasizes stability, efficiency, and high performance. Distinct from earlier methods that rely on token-based importance ratios, GSPO establishes importance ratios via the likelihood of whole sequences, and applies sequence-level techniques for clipping, rewarding, and optimization. Our experiments show that GSPO yields better training efficiency and model outcomes in comparison to the GRPO algorithm, markedly enhances the stability of Mixture-of-Experts (MoE) RL training, and offers promise for streamlining RL system design. GSPO's advantages have played a key role in the substantial progress demonstrated in the latest Qwen3 models.
\end{abstract}

\section{Introduction}
\label{sec:intro}

Reinforcement learning (RL) has become a central methodology for expanding the capabilities of language models \citep{o1,dpsk-r1,qwq32b,qwen3}. By leveraging RL at scale, these models acquire the proficiency required to address complex tasks—such as advanced mathematics and programming—by executing more intricate and prolonged reasoning sequences.

Achieving further scalability in RL, supported by substantial computational resources, relies fundamentally on the preservation of stable and reliable training behavior. Nevertheless, leading RL methods like GRPO \citep{grpo} struggle with pronounced instability when applied to training massive language models, frequently leading to sudden and unrecoverable degradation in performance \citep{qwen3, minimax-m1}. This persistent instability poses a major obstacle to advancing the abilities of language models through ongoing RL optimization.

This study demonstrates that the instability observed in GRPO primarily arises from a fundamental error in how importance sampling weights are applied and validated within its algorithmic framework. This flaw gives rise to training noise with substantial variance, which accumulates as responses grow longer and is exacerbated by the use of clipping, leading ultimately to model breakdown.

To overcome these principal challenges, we introduce \textbf{Group Sequence Policy Optimization (GSPO)}, an innovative reinforcement learning algorithm designed for large language model training. GSPO's primary advancement is its principled definition of the importance ratio, which leverages sequence likelihood as established in previous work \citep{click}, thereby adhering to the foundational tenets of importance sampling. Moreover, GSPO determines normalized rewards by evaluating the advantages across several generated responses for each query, thereby promoting effective alignment between sequence-level reward structures and policy optimization.

The results of our empirical study reveal that GSPO substantially outperforms GRPO in terms of training stability, efficiency, and overall effectiveness. Notably, GSPO addresses the inherent stability issues encountered in the reinforcement learning (RL) training process of extensive Mixture-of-Experts (MoE) architectures, obviating the reliance on elaborate stabilization techniques and highlighting the prospect of streamlining RL frameworks. These strengths of GSPO have played a decisive role in achieving remarkable performance gains in the latest Qwen3 models. We regard GSPO as a dependable and scalable basis for algorithms, supporting further progress in large-scale RL training for language models.

\section{Preliminaries}

\paragraph{Notation}
Throughout this work, an autoregressive language model with parameters $\theta$ is treated as a policy $\pi_\theta$. The symbol $x$ refers to a query, while $\mathcal{D}$ represents the set of queries. For a given query $x$ and its associated response $y$, the probability assigned by the policy $\pi_\theta$ is written as $\pi_\theta (y | x)=\prod_{t=1}^{|y|} \pi_\theta (y_t | x, y_{<t} )$, where $|y|$ stands for the total number of tokens in $y$. Each query-response tuple $(x, y)$ may be evaluated by a verifier $r$, which assigns a reward $r(x, y) \in [ 0, 1 ]$.

\paragraph{Proximal Policy Optimization (PPO)} Using samples generated from the old policy $\pi_{\theta_\text{old}}$, PPO \citep{ppo} constrains the policy update within a proximal region of the old policy through the clipping mechanism. Specifically, PPO employs the following objective for policy optimization (we omit the KL regularization term hereinafter for brevity, as it is not the focus of this paper):

\begin{align}
\mathcal{J}_\text{PPO}(\theta) = \mathbb{E}_{ x \sim \mathcal{D},\, y \sim \pi_{\theta_\text{old}}( \cdot | x) }
\left[ \frac{1}{|y|} \sum_{t=1}^{|y|} 
\min \left( w_{t}(\theta) \widehat{A}_{t},  \, \mathrm{clip} \left( w_{t}(\theta), 1 - {\varepsilon}, 1 + {\varepsilon}\right) \widehat{A}_{t} \right)
\right],
\end{align}

Here, the importance ratio for the token $y_t$ is given by $
w_{t}(\theta) = \frac{ \pi_{\theta} (y_{t} | x, y_{<t}) }{ \pi_{\theta_\text{old}} (y_{t} | x,y_{<t})} $, with the advantage $\widehat{A}_{t}$ corresponding to $y_t$ being inferred using a separate value model, and $\varepsilon$ denotes the threshold for clipping the importance ratios.

A major obstacle in implementing PPO is its significant dependence on the value model. Typically, the value model matches the policy model in scale, which results in substantial demands on both memory and computation. Moreover, the method’s overall performance is directly affected by the precision of the value estimates it produces. Achieving accurate value estimation is already non-trivial, and adapting these estimators for longer outputs and increasingly intricate tasks compounds this difficulty.

\paragraph{Group Relative Policy Optimization (GRPO)} GRPO \citep{grpo} bypasses the need for the value model by computing the relative advantage of each response within a group of responses to the same query. Specifically, GRPO optimizes the following objective:

\begin{align}
\label{equ:grpo}
\mathcal{J}_\text{GRPO}(\theta) = \mathbb{E}_{ x \sim \mathcal{D},\, \{y_i\}_{i=1}^G \sim \pi_{\theta_\text{old}}( \cdot | x) }
\left[ \frac{1}{G} \sum_{i=1}^{G} \frac{1}{|y_i|} \sum_{t=1}^{|y_i|} 
\min \left( w_{i,t}(\theta) \widehat{A}_{i,t},  \, \mathrm{clip} \left( w_{i,t}(\theta), 1 - {\varepsilon}, 1 + {\varepsilon}\right) \widehat{A}_{i,t} \right)
\right],
\end{align}

where $G$ is the number of generated responses to each query $x$ (i.e., the group size), and the importance ratio $w_{i,t}(\theta)$ and advantage $\widehat{A}_{i,t}$ of token $y_{i,t}$ are:

\begin{align}
    w_{i,t}(\theta)=\frac{ \pi_{\theta} (y_{i,t} | x, y_{i,<t}) }{ \pi_{\theta_\text{old}} (y_{i,t} | x,y_{i,<t})},\quad\ 
    \widehat{A}_{i,t} = \widehat{A}_{i} = \frac{r(x, y_i) - \mathrm{mean} \left( \{ r(x, y_i) \}_{i=1}^G \right) }{ \mathrm{std} \left( \{ r(x, y_i) \}_{i=1}^G \right) },
\end{align}

respectively, where all the tokens in $y_i$ share the same advantage as $\widehat{A}_{i}$.

\section{Motivation}
\label{sec:motivation}

As model architectures have increased in parameter count, incorporated greater sparsity (such as Mixture-of-Experts architectures), and generated longer responses, efficient hardware usage during reinforcement learning increasingly requires substantial rollout batch sizes. To address concerns around sample efficiency, a typical strategy divides these extensive rollout batches into several mini-batches that are then used for gradient-based optimization. Such splitting inherently leads to an off-policy regime; in this setting, response sequences $y$ originate from a previous policy $\pi_{\theta_\text{old}}$ instead of the latest policy $\pi_\theta$ currently under optimization. This dynamic underlies the adoption of clipping strategies in algorithms like PPO and GRPO, which serve to filter out samples deemed excessively "off-policy," ensuring that gradient computation remains robust.

While mechanisms like clipping aim to manage this off-policy discrepancy, we identify a more fundamental issue in GRPO: \textit{its objective is ill-posed}. This problem becomes particularly acute when training large models on long-response tasks, leading to catastrophic model collapse. The ill-posed nature of the GRPO objective stems from a misapplication of importance sampling weights. The principle of importance sampling is to estimate the expectation of a function $f$ under a target distribution $\pi_\text{tar}$ by re-weighting samples drawn from a behavior distribution $\pi_\text{beh}$:

\begin{align}
\mathbb{E}_{z \sim \pi_\text{tar}} \left[ f(z) \right] = \mathbb{E}_{z \sim \pi_\text{beh}} \left[ \frac{ \pi_\text{tar} (z) }{ \pi_\text{beh} (z)} \, f(z) \right].
\end{align}

A key aspect here is that effective correction for the discrepancy between distributions requires averaging over a large number of samples ($N \gg 1$) drawn from the behavior policy $\pi_\text{beh}$, so that the importance weighting factor $\frac{ \pi_\text{tar} (z) }{ \pi_\text{beh} (z)}$ performs as intended.

GRPO, in contrast, introduces the importance weight $\frac{ \pi_{\theta} (y_{i,t} | x, y_{i,<t}) }{ \pi_{\theta_\text{old}} (y_{i,t} | x,y_{i,<t})}$ at the level of individual tokens for each time step $t$. Because this weighting relies solely on a single sampled token $y_{i,t}$ from each subsequent-token distribution $\pi_{\theta_\text{old}}( \cdot | x, y_{i, <t})$, it does not properly correct for the underlying distribution gap. Rather, it injects significant variance into gradient updates during training. Such noise can accumulate considerably over longer sequences and is further magnified by the use of clipping. Our experiments have shown that these factors can precipitate catastrophic model collapse—an outcome that is difficult or impossible to recover from. Attempts at resuming training after such collapse, including loading previous checkpoints, adjusting hyperparameters (such as clipping ranges), increasing sequence lengths, or varying RL queries, all fail to restore model performance.

These findings point to a more intrinsic problem with the GRPO methodology. In particular, they highlight the necessity that the unit on which the optimization objective operates should align with the reward's granularity. Given that the reward pertains to whole sequences, implementing off-policy correction on individual tokens is inappropriate. This realization leads us to propose abandoning token-level corrections in favor of approaches that deploy importance sampling and optimize directly at the sequence level.

\section{Algorithm}

\subsection{GSPO: Group Sequence Policy Optimization}

Although the token-level importance weight $\frac{ \pi_{\theta} (y_{i,t} | x, y_{i,<t}) }{ \pi_{\theta_\text{old}} (y_{i,t} | x,y_{i,<t})}$ presents challenges within GRPO, our findings indicate that the \textit{sequence-level} importance weight $\frac{ \pi_{\theta} (y | x) }{ \pi_{\theta_\text{old}} (y | x)}$ possesses a well-defined theoretical significance in language generation tasks. Specifically, this quantity quantifies the extent to which a response $y$ generated by $\pi_{\theta_\text{old}} (\cdot | x)$ diverges from those favored by $\pi_{\theta} (\cdot | x)$, closely corresponding to the sequence-level reward and offering valuable insight for interpreting the function of the clipping procedure.

Based on this straightforward observation, we propose the \textbf{Group Sequence Policy Optimization (GSPO)} algorithm. GSPO employs the following sequence-level optimization objective:

\begin{align}
\mathcal{J}_\text{GSPO} (\theta) =
\mathbb{E}_{ x \sim \mathcal{D},\, \{y_i\}_{i=1}^G \sim \pi_{\theta_\text{old}}( \cdot | x) }
\left[ 
\frac{1}{G} \sum_{i=1}^{G}
\min \left( s_{i}(\theta)  \widehat{A}_{i},  \, \mathrm{clip} \left( s_{i}(\theta), 1 - {\varepsilon}, 1 + {\varepsilon} \right) \widehat{A}_{i} \right) 
\right],
\label{equ:gspo}
\end{align}

where we adopt the group-based advantage estimation:

\begin{align}
\widehat{A}_{i} = \frac{r(x, y_i) - \mathrm{mean} \left( \{ r(x, y_i) \}_{i=1}^G \right) }{ \mathrm{std} \left( \{ r(x, y_i) \}_{i=1}^G \right) },
\end{align}

and define the importance ratio $s_{i}(\theta)$ based on sequence likelihood \citep{click}:

\begin{align}
s_{i}(\theta) = \left( \frac{ \pi_{\theta} (y_i | x) }{ \pi_{\theta_\text{old}} (y_i | x)} \right)^{\frac{1}{|y_i|}}
=
\exp \left( \frac{1}{|y_i|} \sum_{t=1}^{|y_i|} \log \frac{ \pi_{\theta} (y_{i,t} | x, y_{i,<t}) }{ \pi_{\theta_\text{old}} (y_{i,t} | x,y_{i,<t})} \right).
\end{align}

Consequently, GSPO implements clipping at the response level rather than for each token, thereby filtering out excessively ``off-policy'' examples from the gradient calculation. This strategy is consistent with both sequence-level reward assignment and optimization objectives. To ensure that $s_{i}(\theta)$ remains within a consistent numerical interval and to minimize variance, we employ length normalization. In its absence, significant changes in the likelihood of just a few tokens can cause substantial swings in the sequence-level importance ratio, leading to responses with varying lengths necessitating different clipping thresholds. Additionally, we observe that the clipping intervals utilized in GSPO and earlier approaches, such as GRPO, typically vary considerably, mainly as a result of the differing formulations of importance ratios.

\subsection{Gradient Analysis}
\label{subsec:gradient}

We can derive the gradient of the GSPO objective as follows (clipping is omitted for brevity):

\begin{align}
\nabla_{\theta} \mathcal{J}_\text{GSPO} (\theta)
=&\ 
\nabla_{\theta} \mathbb{E}_{ x \sim \mathcal{D},\, \{y_i\}_{i=1}^G \sim \pi_{\theta_\text{old}}( \cdot | x) }
\left[ 
\frac{1}{G} \sum_{i=1}^{G} s_{i}(\theta) \widehat{A}_{i}
\right] \\
= &\ 
\mathbb{E}_{ x \sim \mathcal{D},\, \{y_i\}_{i=1}^G \sim \pi_{\theta_\text{old}}( \cdot | x) }
\left[ 
\frac{1}{G} \sum_{i=1}^{G}
s_{i}(\theta) \widehat{A}_{i}
\cdot \nabla_{\theta} \log s_{i}(\theta)
\right] \\
=&\ 
\mathbb{E}_{ x \sim \mathcal{D},\, \{y_i\}_{i=1}^G \sim \pi_{\theta_\text{old}}( \cdot | x) }
\left[ 
\frac{1}{G} \sum_{i=1}^{G} 
\left( \frac{ \pi_{\theta} (y_i | x) }{ \pi_{\theta_\text{old}} (y_i | x)} \right)^{\frac{1}{|y_i|}} \widehat{A}_{i} 
\cdot \frac{1}{|y_i|} \sum_{t=1}^{|y_i|} \nabla_{\theta} \log \pi_{\theta} (y_{i,t} | x, y_{i,<t}) 
\right].
\label{equ:gspo_grad}
\end{align}

For comparison, the gradient of the GRPO objective is as follows (note that $\widehat{A}_{i,t} = \widehat{A}_{i}$):

\begin{align}
\nabla_{\theta} \mathcal{J}_\text{GRPO} (\theta) 
=&\ 
\nabla_{\theta} \mathbb{E}_{ x \sim \mathcal{D},\, \{y_i\}_{i=1}^G \sim \pi_{\theta_\text{old}}( \cdot | x) }
\left[ \frac{1}{G} \sum_{i=1}^{G} \frac{1}{|y_i|} \sum_{t=1}^{|y_i|} 
w_{i,t}(\theta) \widehat{A}_{i,t}
\right] \\
=&\
\mathbb{E}_{ x \sim \mathcal{D},\, \{y_i\}_{i=1}^G \sim \pi_{\theta_\text{old}}( \cdot | x) }
\left[ \frac{1}{G} \sum_{i=1}^{G} \widehat{A}_{i} 
\cdot \frac{1}{|y_i|} \sum_{t=1}^{|y_i|} 
\frac{ \pi_{\theta} (y_{i,t} | x, y_{i,<t}) }{ \pi_{\theta_\text{old}} (y_{i,t} | x,y_{i,<t})} 
\nabla_{\theta} \log \pi_{\theta} (y_{i,t} | x, y_{i,<t})  
\right].
\end{align}

The key difference between GSPO and GRPO centers on the manner in which they assign weights to the gradients of token log likelihoods. GRPO utilizes ``importance weights'' for each token, calculated as $\frac{ \pi_{\theta} (y_{i,t} | x, y_{i,<t}) }{ \pi_{\theta_\text{old}} (y_{i,t} | x,y_{i,<t})}$, which means the gradient contributions are scaled unevenly depending on these values. These weights can span the intervals $(0, 1+\varepsilon]$ when $\widehat{A}_{i} >0$, or $[1-\varepsilon, +\infty)$ when $\widehat{A}_{i} < 0$, making them significant and potentially problematic, as their effects can accumulate during training, introducing volatility and unpredictability. On the other hand, GSPO treats all tokens within a generated sequence with uniform weights, thereby removing this source of instability inherent to GRPO.

\subsection{GSPO-token: A Token-level Objective Variant}

In scenarios like multi-turn RL, we may desire a finer-grained advantage adjustment than the sequence level. To this end, we introduce a token-level objective variant of GSPO, namely \textbf{GSPO-token}, to allow token-wise advantage customization:

\begin{align}
\mathcal{J}_\text{GSPO-token}(\theta)
=
\mathbb{E}_{ x \sim \mathcal{D},\, \{y_i\}_{i=1}^G \sim \pi_{\theta_\text{old}}( \cdot | x) }
\left[ 
\frac{1}{G} \sum_{i=1}^{G} \frac{1}{|y_i|} \sum_{t=1}^{|y_i|} 
\min \left( s_{i,t}(\theta) \widehat{A}_{i,t},  \, \mathrm{clip} \left( s_{i,t}(\theta), 1 - {\varepsilon}, 1 + {\varepsilon} \right) \widehat{A}_{i,t} \right) 
\right],
\label{equ:gspo-token}
\end{align}

where

\begin{align}
s_{i,t}(\theta) 
= \mathrm{sg} \left[ s_{i}(\theta) \right]  \cdot \frac{ \pi_{\theta} (y_{i,t} | x, y_{i,<t}) }{ \mathrm{sg} \left[ \pi_{\theta} (y_{i,t} | x, y_{i,<t}) \right] },
\end{align}

and $\mathrm{sg}[\cdot]$ denotes only taking the numerical value but stopping the gradient, corresponding to the \texttt{detach} operation in PyTorch. The gradient of GSPO-token can be derived as:

\begin{align}
\nabla_{\theta} \mathcal{J}_\text{GSPO-token}(\theta)
=&\ 
\nabla_{\theta} \mathbb{E}_{ x \sim \mathcal{D},\, \{y_i\}_{i=1}^G \sim \pi_{\theta_\text{old}}( \cdot | x) }
\left[ 
\frac{1}{G} \sum_{i=1}^{G}
\frac{1}{|y_i|} \sum_{t=1}^{|y_i|} 
s_{i,t}(\theta) \widehat{A}_{i,t}
\right] \\
=&\ 
\mathbb{E}_{ x \sim \mathcal{D},\, \{y_i\}_{i=1}^G \sim \pi_{\theta_\text{old}}( \cdot | x) }
\left[ 
\frac{1}{G} \sum_{i=1}^{G} s_i (\theta) \cdot \frac{1}{|y_i|} \sum_{t=1}^{|y_i|} 
\widehat{A}_{i,t} \frac{ \nabla_{\theta} \pi_{\theta} (y_{i,t} | x, y_{i,<t}) }{ \pi_{\theta} (y_{i,t} | x, y_{i,<t}) }
\right] \\
=&\ 
\mathbb{E}_{ x \sim \mathcal{D},\, \{y_i\}_{i=1}^G \sim \pi_{\theta_\text{old}}( \cdot | x) }
\left[ 
\frac{1}{G} \sum_{i=1}^{G}  \left( \frac{ \pi_{\theta} (y_i | x) }{ \pi_{\theta_\text{old}} (y_i | x)} \right)^{\frac{1}{|y_i|}}
\cdot \frac{1}{|y_i|} \sum_{t=1}^{|y_i|}
\widehat{A}_{i,t} \nabla_{\theta} \log \pi_{\theta} (y_{i,t} | x, y_{i,<t})  
\right].
\label{equ:gspo-token_grad}
\end{align}

It should be emphasized that the expression $\frac{ \pi_{\theta} (y_{i,t} | x, y_{i,<t}) }{ \mathrm{sg} \left[ \pi_{\theta} (y_{i,t} | x, y_{i,<t}) \right] }$ evaluates to 1, meaning that $s_{i,t}(\theta)$ effectively yields the same numerical result as $s_{i}(\theta)$. When examining Equation~\eqref{equ:gspo} and~\eqref{equ:gspo-token} as well as Equation~\eqref{equ:gspo_grad} versus~\eqref{equ:gspo-token_grad}, GSPO-token matches GSPO numerically with respect to their optimization targets, clipping criteria, and analytical gradients, provided all token advantages within the output $y_i$ are identical (specifically, if $\widehat{A}_{i,t} = \widehat{A}_{i}$). Nevertheless, GSPO-token allows for greater adaptability, enabling variation in advantage assignments at the token level.

\section{Experiments and Discussion}

\subsection{Empirical Results}

Our experiments utilize a cold-start model initialized from Qwen3-30B-A3B-Base, and we present both training reward trajectories and evaluation results on the AIME'24 (average Pass@1 from 32 samples), LiveCodeBench (202410-202502, average Pass@1 from 8 samples), and CodeForces (Elo Rating) benchmarks. Throughout RL training, each batch of rollout samples is split into four mini-batches to perform gradient updates. For GSPO, the clipping intervals in Equation~\eqref{equ:gspo} are set to 3e-4 and 4e-4 for the lower and upper bounds, respectively. We benchmark this setup against GRPO, designating the clipping limits in Equation~\eqref{equ:grpo} as 0.2 and 0.27, following meticulous tuning to allow an equitable comparison. It is important to mention that GRPO requires the use of the Routing Replay training technique for stable convergence in MoE RL scenarios, which we elaborate on in \S~\ref{subsec:routing_replay}, whereas \textit{GSPO eliminates this requirement}.

As illustrated in Figure~\ref{fig:results}, training utilizing GSPO exhibits consistent stability. Our findings indicate that \textbf{GSPO enables sustained gains in performance by augmenting computational resources, systematically refreshing the query set, and lengthening generation outputs}. Additionally, GSPO outperforms GRPO in terms of training efficiency, attaining higher accuracy and improved benchmark results when matched for compute and number of queries used. Lastly, GSPO has been effectively integrated into RL training for the most recent Qwen3 models, robustly validating GSPO’s capability to enhance RL scaling in large language model contexts.

\subsection{Curious Observation on Clipping Fractions}

An important differentiator between GSPO and GRPO lies in GSPO's approach of clipping at the response level rather than at the token level. As illustrated in Figure~\ref{fig:clipping}, this methodology results in a two orders of magnitude difference in the proportion of tokens clipped, with GSPO removing a far greater fraction than GRPO (a pattern that persists even when the clipping ranges are adjusted). Notably, although GSPO clips many more tokens—thus utilizing fewer for model training or gradient computation—it nonetheless demonstrates greater training efficiency compared to GRPO. This somewhat surprising outcome, where extensive token clipping corresponds to improved efficiency, suggests that gradient estimates at the token level in GRPO are intrinsically noisy and suboptimal for exploiting samples. Conversely, GSPO’s reliance on sequence-level gradients provides a more dependable and productive signal for learning.

\subsection{Benefit of GSPO for MoE Training}
\label{subsec:routing_replay}

\paragraph{Background}
In contrast with RL training of traditional dense models, MoE models pose distinct stability issues due to their inherently sparse activation pattern. Specifically, when utilizing the GRPO algorithm, the phenomenon we term \textit{expert-activation volatility} in MoE architectures poses a barrier to successful RL convergence. More precisely, we observed that following one or more applications of gradient descent, the subset of experts engaged to produce identical outputs may change drastically. As an illustrative case, for the 48-layer Qwen3-30B-A3B-Base model, following each RL gradient update, nearly 10\% of the experts triggered for a given rollout sample by the updated policy $\pi_{\theta}$ are different from those activated under the preceding policy $\pi_{\theta_\text{old}}$.

This behavior intensifies with increasing model depth in MoE systems, leading to pronounced oscillations in the token-level importance ratios $w_{i,t}(\theta) = \frac{ \pi_{\theta} (y_{i,t} | x, y_{i,<t}) }{ \pi_{\theta_\text{old}} (y_{i,t} | x,y_{i,<t})}$ and, as outlined in \S~\ref{sec:motivation} and~\ref{subsec:gradient}, these ratios become unreliable. Such instability ultimately interferes with normal RL training convergence.

\paragraph{Our Previous Approach} In earlier work, we adopted the \textbf{Routing Replay} training strategy to address this issue. Our method involves storing the set of experts activated by $\pi_{\theta_\text{old}}$ and subsequently ``replaying'' these expert configurations during the evaluation of $\pi_{\theta}$ as we calculate importance ratios, $w_{i,t}(\theta) = \frac{ \pi_{\theta} (y_{i,t} | x, y_{i,<t}) }{ \pi_{\theta_\text{old}} (y_{i,t} | x,y_{i,<t})}$. By ensuring that, for each token $y_{i,t}$, both $\pi_{\theta} (y_{i,t} | x, y_{i,<t})$ and $\pi_{\theta_\text{old}} (y_{i,t} | x, y_{i,<t})$ use an identical expert routing, the approach brings consistency to the computation of token-level importance ratios and facilitates stable optimization for the same network subset across gradient updates. As depicted in Figure~\ref{fig:routing_replay}, Routing Replay is a crucial component for guaranteeing stable convergence in GRPO training processes for MoE models.

\paragraph{Benefit of GSPO} While Routing Replay allows MoE models to converge under GRPO training, it introduces extra costs in memory and communication and may restrict the effective utilization of the model's capacity. GSPO, as depicted in Figure~\ref{fig:results}, removes the reliance on Routing Replay, readily computes the importance ratios $s_i(\theta)$ in the standard manner, and ensures normal convergence with consistent optimization. The central realization is that GSPO prioritizes the overall sequence likelihood (i.e., $\pi_{\theta} (y_i | x)$), remaining unaffected by fluctuations in individual token likelihoods (i.e., $\pi_{\theta} (y_{i,t} | x, y_{i,<t})$). Given that the MoE architecture retains robust language modeling abilities, the sequence likelihood remains stable across training. Thus, GSPO inherently mitigates the instability in expert activation present in MoE models, eliminating the necessity for workaround solutions such as Routing Replay. This leads to both a more straightforward and reliable training process and enables full utilization of model capacity without imposed limitations.

\subsection{Benefit of GSPO for RL Infrastructure}

Because training engines such as Megatron and inference engines like SGLang and vLLM exhibit differences in precision, it is common practice to recalculate the likelihoods of generated responses under the previous policy $\pi_{\theta_\text{old}}$ using the training engine. In contrast, GSPO optimizes at the sequence level rather than the token level, which inherently provides greater robustness to precision mismatches. As a result, GSPO enables the direct utilization of likelihoods produced by the inference engine during optimization, eliminating the necessity for recomputation with the training engine. This approach offers particular advantages in contexts like partial rollouts, multi-turn RL, and disaggregated architectures where training and inference are separated.

\section{Conclusion}

We introduce Group Sequence Policy Optimization (GSPO), a novel reinforcement learning approach tailored for the training of large language models. GSPO leverages the principle of importance sampling by computing importance ratios using sequence likelihood, and implements mechanisms such as sequence-level clipping, reward assignment, and optimization. Empirical results show that GSPO achieves substantially greater training stability, efficiency, and overall performance relative to GRPO, and is especially effective for large-scale RL training of MoE models. This advancement has enabled the significant progress observed in the latest Qwen3 models. With GSPO as a robust and scalable foundation, we aim to further extend the capabilities of RL at scale, anticipating essential breakthroughs in machine intelligence.

\end{document}